\thispagestyle{empty}
\begin{tcolorbox}[enhanced, colback=hblue, colframe=hblue, boxrule=0pt, arc=4pt, left=16pt, right=16pt, top=22pt, bottom=22pt]
  \color{white}\sffamily
  {\large\bfseries Bewijzen in de Wiskunde \textperiodcentered\ 2026--2027}\\[10pt]
  {\fontsize{34}{38}\selectfont\bfseries Workbook}\\[6pt]
  {\fontsize{20}{24}\selectfont Chapters 0, 1 and 2 of the book, with every exercise worked}\\[14pt]
  {\normalsize Minimal theory, taken from the book's own definitions and strategies, followed immediately by the book's own exercises.\\
   Every exercise in the covered pages is either worked in the text or solved in Appendix S.}
\end{tcolorbox}

\vspace{10pt}
\begin{tcolorbox}[enhanced, colback=white, colframe=hblue2, boxrule=0.8pt, arc=2pt, left=10pt, right=10pt, top=8pt, bottom=8pt]
\sffamily\small\sloppy
\begin{itemize}[leftmargin=1.2em]
\item \textbf{\textcolor{hblue}{Source.}} C.~Newstead, \emph{An Infinite Descent into Pure Mathematics}, Utrecht adaptation (version of 29 July 2026). All numbers (Definition 0.36, Strategy 1.1.28, Exercise 2.2.26, \dots) are the book's numbers, so you can look everything up.
\item \textbf{\textcolor{hblue}{Covered.}} Chapter 0 (pages 1--25, including 0.E), Chapter 1 (pages 29--75, including 1.E), Section 2.1 (pages 78--84) and Section 2.2 up to Definition 2.2.46 (pages 85--93).
\item \textbf{\textcolor{hblue}{Not covered}} (not in the uploaded pages): the rest of \S2.2 from tuples and Cartesian products onwards, the chapter exercises 2.E, and Chapters 3--5. Upload those pages and this workbook can be extended the same way.
\end{itemize}
\end{tcolorbox}

\vspace{8pt}
\hsub{How this workbook works}
Every section is a sequence of short cycles: \textbf{one piece of theory, then the exercises that use it.}

\begin{center}\sffamily\small
\begin{tabular}{@{}p{0.27\linewidth}p{0.68\linewidth}@{}}
\colorbox{hyellow}{\strut\textbf{Definition n.m}} & The book's definition, condensed. Learn it word for word: every proof starts by unfolding one of these.\\[3pt]
\colorbox{hlight}{\strut\textbf{Theorem / Proposition / Axiom}} & A result you may quote by number in a proof.\\[3pt]
\colorbox{hgreen}{\strut\textbf{Strategy n.m}} & The book's proof strategies. Each one is a rule for what to write next.\\[3pt]
\colorbox{white}{\strut\textcolor{hblue}{\textbf{Exercise n.m}} \textcolor{horange}{\textbf{badge}}} & A book exercise worked in full, in the style the book uses. Cover the solution and try first.\\[3pt]
\colorbox{hpurplebg}{\strut\textbf{NOW YOU}} & The remaining book exercises for that piece of theory. Do them; every one is solved in Appendix S.\\[3pt]
\colorbox{hredbg}{\strut\textbf{Watch out}} & The book's own warnings (``Common error'', asides) and the places where marks are lost.
\end{tabular}
\end{center}

\begin{key}[{Flags from the course planning}]
\fE\ = \textbf{essential} exercise in the planning \quad \fR\ = \textbf{recommended} \quad \fO\ = in the book but not in the planning (the polynomial and complex-number part of Chapter 0, and the questions that use it, are outside the course; they are still solved here so that you can do \emph{any} question in the book).\\[3pt]
Quiz 1 (23 Sep): numbers and \S1.1. \quad Quiz 2 (7 Oct): sets and functions. \quad Exam 4 Nov.
\end{key}

\begin{key}[{Conventions of this book}]
\begin{itemize}
\item $\N=\{0,1,2,\dots\}$: $0$ is a natural number (Definition 0.6). $[n]=\{k\in\N\mid k<n\}$ (Definition 0.25), so $[0]=\varnothing$ and $[3]=\{0,1,2\}$.
\item Truth values in truth tables are $\checkmark$ (true) and $\times$ (false), as in the book.
\item ``$b$ divides $a$'' means $a=qb$ for some $q\in\Z$ (Definition 0.36); it is about multiplication, never about the fraction $a/b$.
\item ``Odd'' means ``not even'' (Definition 0.41); that $a$ is odd iff $a=2b+1$ is a \emph{proposition} (1.1.58), not the definition.
\item $\square$ ends a proof. \reason{grey text in braces} names the definition, strategy or rule used for a step.
\end{itemize}
\end{key}
